\section{Methodology}

\subsection{System Overview}
Our system leverages the PaddleOCR framework to extract structured information from receipt images. The workflow follows a tiered architecture where images are first localized for text regions and then transcribed into digital characters.


\subsection{Image Processing}

Before we can extract text from receipt images using OCR, we need to clean up and enhance the images first. Receipt photos taken with phones often have problems like poor lighting, blur, noise, and low contrast. Our preprocessing pipeline solves these issues step-by-step to prepare clean images for the OCR engine.

\subsubsection{Converting to Grayscale}

The first step is converting color images to grayscale. However, not all receipts are the same - some use black ink while others use blue ink (common in printed invoices). Regular grayscale conversion doesn't work well for blue-ink receipts because it reduces the contrast between text and background.

To handle this, we first check if a receipt has blue ink by comparing the blue and red color channels. If the average blue intensity is more than 10 units higher than the red intensity, we classify it as a blue-ink receipt.

For blue-ink receipts, we use a special conversion formula that ignores the blue channel and only combines the red and green channels:
\begin{equation}
I_{\text{gray}} = 0.5 \times R + 0.5 \times G
\end{equation}

This gives us better text contrast. For normal black-ink receipts, we just use the standard OpenCV grayscale conversion.

\subsection{Text Detection}
\textbf{Purpose:} This component is responsible for locating and generating bounding boxes around text areas, including straight, curved, or rotated text.

\subsubsection{DBNet++ (Differentiable Binarization)}
\textbf{Overview:} A segmentation-based approach highly effective for locating text regardless of orientation or irregular shapes, utilized as the primary detection module in PP-OCRv5.

\textbf{Architecture (PP-OCRv5):}
\begin{itemize}
    \item \textbf{Backbone:} Employs high-performance architectures like \textbf{PP-LCNet} or \textbf{MobileNetV3} specifically optimized for extreme inference speeds on mobile and edge devices, alongside larger networks like \textbf{ResNet-50} for server-side scenarios.
    \item \textbf{Neck (RSE-FPN):} Uses an advanced \textbf{Feature Pyramid Network (FPN)} with \textbf{RSE-FPN} (Residual Squeeze-and-Excitation FPN) modules to fuse multi-scale features and preserve spatial information.
    \item \textbf{Detection Head:} Processes output from the Neck through convolutional and de-convolutional layers to generate:
    \begin{itemize}
        \item \textbf{Probability map:} Indicates text presence at the pixel level.
        \item \textbf{Threshold map:} Dynamically calculates binarization thresholds for varied lighting and contrast conditions.
    \end{itemize}
\end{itemize}

\textbf{Processing and Optimization:}
\begin{itemize}
    \item \textbf{Differentiable Binarization (DB):} Produces a Binary Map to extract sharp, polygon-based text boundaries, optimized using binary cross-entropy (BCE) loss.
    \item \textbf{Geometric Refinement:} To refine text localization and bounding box predictions, core geometric mechanisms are applied:
    \begin{itemize}
        \item \textbf{IoU (Intersection over Union):} Measures the overlap between the predicted box and the ground truth.
        \item \textbf{NMS (Non-Maximum Suppression):} Removes redundant overlapping detections, retaining only the box with the highest confidence score for each text instance.
    \end{itemize}
\end{itemize}

\subsubsection{EAST (Efficient and Accurate Scene Text detector)}
\textbf{Overview:} A U-Net style structure designed for fast and accurate scene text localization, bypassing the need for computationally expensive candidates or word grouping.

\textbf{Architecture:}
\begin{itemize}
    \item \textbf{Feature Extractor:} Utilizes a \textbf{VGG16} backbone (with Batch Normalization) to extract multi-level feature maps from the input image.
    \item \textbf{Feature Merger:} Recovers the spatial resolution of features through non-local merging of top-down and bottom-up layers using bilinear interpolation and concatenation.
    \item \textbf{Output Layer:} Predicts two branches: a \textbf{Score map} (probability of text presence) and a \textbf{Geometry map} (\textbf{RBOX format} including 4 distance offsets and 1 rotation angle).
\end{itemize}

\textbf{Processing and Optimization:}
\begin{itemize}
    \item \textbf{Joint Loss:} Combines \textbf{Dice Loss} (for score map imbalance) with \textbf{Log-IoU} (for bounding box dimensions) and a \textbf{Cosine Angle Loss} (for rotation angle periodicity).
    \item \textbf{Restoration and LANMS:} Decodes the RBOX geometry back into 4-point quadrilaterals followed by \textbf{Locality-Aware NMS (LANMS)} to merge neighboring quadrilaterals from the same text instance.
\end{itemize}

\subsection{Text Recognition}
\textbf{Purpose:} Converts cropped text regions into digital character strings.

\subsubsection{CRNN (Convolutional Recurrent Neural Network)}
\textbf{Overview:} A hybrid architecture that combines CNN-based feature extraction with RNN-based sequence modeling.

\textbf{Architecture:}
\begin{itemize}
    \item \textbf{Convolutional Stage:} A customized \textbf{VGG16-based extractor} with 7 convolutional blocks. Later MaxPool layers utilize a rectangular kernel and stride of \textbf{$(2, 1)$} instead of $(2, 2)$ to preserve horizontal spatial resolution for character sequences.
    \item \textbf{Recurrent Stage:} A deep \textbf{Bidirectional LSTM (Bi-LSTM)} with two layers and 256 hidden units to model contextual dependencies between sequential features.
    \item \textbf{Transcription Layer:} A final linear layer mapping RNN outputs to the target character space.
\end{itemize}

\textbf{Processing and Optimization:}
\begin{itemize}
    \item \textbf{CTC Loss:} Trained using \textbf{Connectionist Temporal Classification} to handle variable-length sequences without requiring precise per-character alignment.
    \item \textbf{Decoding:} During inference, the model utilizes a \textbf{CTC Greedy Search} algorithm to handle blank tokens and consolidate repeated characters for accurate text reconstruction.
\end{itemize}

\subsubsection{SVTR (Single Visual Model for Scene Text Recognition)}
\textbf{Overview:} A specialized transformer-based model for text recognition in PP-OCRv5 that completely replaces Recurrent layers with visual self-attention. To balance high recognition accuracy with low real-world latency, the model adopts an \textbf{SVTR-LCNet} hybrid structure, combining lightweight CNN blocks (PP-LCNet) in the initial stages with SVTR transformer blocks in the final stage.

\textbf{Architecture (PP-OCRv5):}
\begin{enumerate}
    \item \textbf{Progressive Overlapping Patch Embedding:} Converts images into "character components" using two consecutive \textbf{$3 \times 3$ convolutional layers (stride 2)} with Batch Normalization, preserving spatial continuity at character edges.
    \item \textbf{Mixing Blocks:} Interleaves \textbf{Local Mixing} (using a \textbf{$7 \times 11$ sliding window} for stroke morphology) and \textbf{Global Mixing} (capturing inter-character context across the entire image). Local blocks are typically sequenced before Global blocks.
\end{enumerate}

\textbf{Processing and Optimization:}
\begin{itemize}
    \item \textbf{Structural Operators:} Uses \textbf{Merging} operators (\textbf{$3 \times 3$ convolution, stride $(2, 1)$}) to reduce feature map height while preserving width, followed by \textbf{Combining} operators to form a final 1D feature sequence.
    \item \textbf{Prediction:} Features are decoded via a Fully Connected layer combined with \textbf{CTC loss} to resolve the final character string.
\end{itemize}
